\section{执行摘要}
\label{sec:executive-summary}

2024--2026 年，量子数据用于经典机器学习的产业与研究格局呈现\textbf{benchmark 降温、场景分化、混合范式成熟}三大特征。

\subsection{核心结论}

\begin{enumerate}[leftmargin=*]
  \item \textbf{通用 QML 未证明显著优于经典}：Bowles 等~\cite{Bowles2024BetterThanClassical}、Schnabel \& Roth~\cite{Schnabel2025QKMBenchmark}、Fellner 等~\cite{Fellner2026QMLTimeSeries} 的大规模对照表明，在小数据集分类与时序任务上，开箱即用量子模型整体不优于经典基线。
  \item \textbf{路径分化明确}：
  \begin{itemize}
    \item \pathone{}（标签一一对应）：处理器代理~\cite{PredictiveSurrogates2026}、ML-QEM~\cite{SZNE2025,NoiseAgnosticQEM2025}、VQE warm-start~\cite{Qracle2025} 已具工程价值。
    \item \pathtwo{}（原生测量数据）：Cho \& Kim~\cite{Cho2024ExpMLQuantumData}、ShadowGPT~\cite{ShadowGPT2024}、LLM4QPE~\cite{LLM4QPE2024} 证明 foundation 范式可行。
  \item \textbf{产业仍处 pilot 阶段}：Quantum AI 细分市场 2025 约 4 亿美元，2030 约 18 亿美元（CAGR $\sim$34\%），企业以 hybrid pilot 为主~\cite{EUQuantumAIWhitePaper2025,QEDCQuantumAI2025}。
\end{enumerate}

\subsection{调研范围与方法}

\begin{itemize}[leftmargin=*]
  \item \textbf{时间窗口}：2024-07 至 2026-07
  \item \textbf{文献规模}：83 篇代表性 BibTeX 条目（见 \texttt{quantum\_data\_ml\_survey.bib}）
  \item \textbf{分类框架}：路径 1 / 路径 2 / 混合 + Benchmark / 理论 / 行业
\end{itemize}
